% Ported from an earlier report section (one-time conversion; edit here).
\subsection{The effort confound \texorpdfstring{\screen}{[screen]}}\label{vt:sec:effort}

An external reviewer asked the question this campaign should have asked itself: when the sticky router chose Sonnet,
was its session the same as the plain-Sonnet control? It was not. Every routed Sonnet main request ran at
\emph{medium} reasoning effort (\RvEffMediumSticky\ of \RvEffStickyTotal\ requests in the sticky sessions that chose
Sonnet), because the bundle's composed configuration sets that start effort. The plain-Sonnet control sets no effort
(\RvEffNoneControl\ of \RvEffControlTotal\ requests), which the design document describes as the provider default,
high. So the comparison ``routed versus plain Sonnet'' also compares two effort levels.

\Cref{vt:tab:effort} gives the full numbers for this part.

\begin{table}[tbp]
\centering\small
\caption{Sticky sessions that chose Sonnet against the plain-Sonnet control on the same scenario and repetition
(geometric-mean cost ratio, 95\,\% scenario-cluster bootstrap intervals). Exploratory.}
\label{vt:tab:effort}
\begin{tabular*}{\textwidth}{@{\extracolsep{\fill}}l r r r r@{}}
\toprule
 & \multicolumn{2}{c}{All scenarios} & \multicolumn{2}{c}{Test split} \\
\cmidrule(lr){2-3}\cmidrule(l){4-5}
Host & sessions & cost ratio & sessions & cost ratio \\
\midrule
Fable & \RvEffFableAllN & \RvEffFableAll\ (\RvEffFableAllCI) & \RvEffFableTestN & \RvEffFableTest\ (\RvEffFableTestCI) \\
Opus & \RvEffOpusAllN & \RvEffOpusAll\ (\RvEffOpusAllCI) & \RvEffOpusTestN & \RvEffOpusTest\ (\RvEffOpusTestCI) \\
\bottomrule
\end{tabular*}
\end{table}

\paragraph{What the medium-effort sessions did differently} Against the control on the same scenarios, the sticky
sessions that ran on Sonnet made about \RvEffReqFable\,\% fewer requests on Fable (\RvEffReqOpus\,\% on Opus),
produced about \RvEffOutFable\,\% fewer output tokens (\RvEffOutOpus\,\%) and finished about \RvEffTimeFable\,\%
faster (\RvEffTimeOpus\,\%). Quality showed no detectable difference: turn-pass fraction moved by \RvEffDtpFable\
on Fable (95\,\% CI \RvEffDtpFableCI) and \RvEffDtpOpus\ on Opus (\RvEffDtpOpusCI). The sessions were not using the
read shortcut; the difference is the effort setting, possibly together with the bundle's own context, which this
data cannot separate.

\paragraph{Two further corrections}
\begin{itemize}
  \item \textbf{The control ran in the Opus waves only.} Because it does not depend on the host, the plain-Sonnet
    control ran once per scenario and repetition, alongside the Opus sessions. Fable sessions and their Sonnet
    control started a median \RvSonnetWaveMedianH\ hours apart (at most \RvSonnetWaveMaxH\,h). The Opus result,
    with the control in the same wave, is about the same size, so timing is not what drives the gap.
  \item \textbf{Sticky sessions that kept the host cost more than the plain host.} The \RvStickyHostNFable\ sticky
    sessions per host that stayed on the host model cost \RvStickyHostFable\X\ the plain host on Fable and
    \RvStickyHostOpus\X\ on Opus. On Fable, sticky-on-Sonnet cost \RvEffVsAnchorFable\X\ the plain host where the
    control cost \RvCtlVsAnchorFable\X\ on the same scenarios; the dearer host-kept sessions offset most of that
    difference, which is why sticky and plain Sonnet look alike overall. Because those host-kept sessions ran the
    host near its default effort mix, \RvStickyHostFable\X\ and \RvStickyHostOpus\X\ are a near-effort-matched estimate
    of the bundle's own overhead on the host, about \RgOverheadFable--\RgOverheadOpus\,\%. With \RvStickyHostNFable\
    sessions per host and an approximate effort match, this is an indication, not a finding.
\end{itemize}

\subsubsection{Follow-up: plain Sonnet at medium versus default effort \texorpdfstring{[preregistered, provenance-limited]}{[preregistered, provenance-limited]}}\label{vt:sec:effortfollow}

To separate the effort setting from the routing, a preregistered, provenance-limited follow-up (effort-control-v1) ran plain Sonnet
twice on the \EcScen\ test scenarios, \EcReps\ repetitions each: once at the provider's default effort and once at
medium effort, nothing else changed (all \RgEcMainMedium\ of \RgEcMainTotal\ main requests of the medium arm ran at
medium; its \RgEcNonMainDefault\ default-effort requests are all background, not main-loop). Each pair started together from the same frozen workspace with separate cache
nonces, so this comparison is concurrent. Intervals come from \EcResamples\ scenario-cluster resamples (seed
\EcSeed). The preregistered hypothesis (medium is cheaper, with non-inferior turn-pass fraction) was \EcHE.

\Cref{vt:tab:effortfollow} gives the full numbers for this part.

\begin{table}[tbp]
\centering\small
\caption{Effort-control follow-up (evidence package in \cref{tab:evidence}). Cost rows are
geometric-mean ratios of tools-normalized session cost; below 1 means the first arm is cheaper.}
\label{vt:tab:effortfollow}
\begin{tabular*}{\textwidth}{@{\extracolsep{\fill}}>{\raggedright\arraybackslash}p{0.5\textwidth} r r r@{}}
\toprule
Comparison & n & estimate & 95\,\% CI \\
\midrule
Medium / default effort, plain Sonnet (confirmatory, concurrent) & 46 pairs & 0.821 & 0.781--0.861 \\
Medium plain Sonnet / main-v1 sticky-on-Sonnet (exploratory, not concurrent) & 22 scenarios & 0.958 & 0.913--0.999 \\
\midrule
Turn-pass difference, medium minus default (unfiltered) & 46 pairs & $-$0.003 & $-$0.026 to 0.021 \\
\bottomrule
\end{tabular*}

\end{table}

\paragraph{Result} Medium effort cost \EcRatio\X\ the default (95\,\% CI \EcCI), about \EcCutPct\,\% cheaper, and
was cheaper in \EcCheaperPairs\ of \EcPairs\ pairs. Turn-pass fraction moved by \EcQual\ (\EcQualCI): \EcQualVerdict\
at the \NIMarginPts-point margin. An exploratory comparison of medium plain Sonnet with the main-campaign sticky
sessions that ran on Sonnet (also at medium effort) gives \EcSecRatio\ (\EcSecCI, \EcSecN\ scenarios). Those
sessions ran on different dates, so this second comparison is not concurrent and has no decision rule. Read the
other way round, sticky-on-Sonnet cost about \EcSecGapPct\,\% more than medium plain Sonnet; that gap is within what
bundle behaviour or timing could produce.

\paragraph{Provenance caveat} The follow-up's schedule and plan were created \EcScheduleBefore\ \emph{before} the
preregistration was committed (\code{\EcPreregSha}), and the first analysed session started \EcFirstAfter\ after it.
The design file and the analysis script were first committed in that same commit, so the evidence package cannot
prove that the working-tree copies used to build the schedule match the committed ones. No outcome data from the
analysed sessions existed before the commit. The follow-up cost \$\EcSpend\ (including a \$\EcPreflight\ preflight).

\begin{keyidea}
The effort setting alone explains most of the gap between sticky-on-Sonnet and plain Sonnet: in a concurrent test,
medium effort cost \EcRatio\X\ the default. On Fable~5.1 at the \cref{vt:tab:prices} prices, the saving comes from
running the session on Sonnet at medium effort. On the Sonnet sessions the router's extra machinery did not help: sticky-on-Sonnet cost about \EcSecGapPct\,\% more
than plain Sonnet at medium effort (\EcSecRatio\ the other way round; exploratory, not concurrent), and the
\RvStickyHostNFable\ sessions per host that kept the host cost \RvStickyHostFableTwo\X\ (Fable) and
\RvStickyHostOpusTwo\X\ (Opus) the plain host. This campaign shows no benefit of this router beyond the decision to use
Sonnet. The finding also applies to anyone running plain Sonnet:
medium effort cut cost by about \EcCutPct\,\% with no detectable turn-pass loss on these \EcScen\ scenarios (\RvDocsTest\ docs, \RvExplainTest\ explain, \RvReviewTest\ review); the main campaign
hints at a loss on docs, explain and review work (sticky-on-Sonnet at medium minus the default-effort control:
\RvKTDtp\ [\RvKTDtpCI] on those \RvKTN\ scenarios against \RvRestDtp\ [\RvRestDtpCI] on the rest; post hoc, and
confounded with the bundle's context) that the follow-up is too small to settle.
\end{keyidea}

\subsubsection{Follow-up: plain Fable at medium versus default effort \texorpdfstring{[preregistered]}{[preregistered]}}\label{vt:sec:effortfable}

The round-2 review asked how the host itself responds to effort. A second preregistered follow-up
(effort-control-fable-v1) ran plain Fable~5.1 twice on the \EcScen\ test scenarios, \EcReps\ repetitions each, at the
provider default and at medium effort, concurrently, with the same design and \EcResamples\ scenario-cluster resamples
(seed \EcSeed). This time the preregistration was committed (\code{\EfPreregSha}) before the schedule was created
(\EfScheduleAfter\ later); a throwaway feasibility schedule and a two-session mechanism check made before the commit
are disclosed in its provenance record.

\Cref{vt:tab:effortfable} gives the full numbers for this part.

\begin{table}[tbp]
\centering\small
\caption{Fable effort follow-up (evidence package in \cref{tab:evidence}). Cost rows are
geometric-mean ratios of tools-normalized session cost; below 1 means the first arm is cheaper.}
\label{vt:tab:effortfable}
{\footnotesize\setlength{\tabcolsep}{2pt}\begin{tabular*}{\textwidth}{@{\extracolsep{\fill}}>{\raggedright\arraybackslash}p{0.5\textwidth} r r r@{}}
\toprule
Comparison & n & estimate & 95\,\% CI \\
\midrule
Fable medium / default effort (confirmatory, concurrent) & 46 pairs (45 cost-valid) & 0.860 & 0.833--0.885 \\
Same, excluding the audit-flagged pair (sensitivity) & 44 pairs & 0.854 & 0.825--0.882 \\
Fable medium / main-campaign Fable sticky (exploratory, not concurrent) & 23 scenarios & 1.546 & 1.343--1.759 \\
Fable medium / main-campaign Fable anchor (exploratory, not concurrent) & 23 scenarios & 0.829 & 0.790--0.866 \\
\midrule
Turn-pass difference, medium minus default (unfiltered) & 46 pairs & 0.033 & $-$0.002 to 0.080 \\
\bottomrule
\end{tabular*}
}
\end{table}

\paragraph{Result} Plain Fable at medium effort cost \EfRatio\X\ the default (95\,\% CI \EfCI; \EfValid\ cost-valid of
\EfPairs\ pairs), about \EfCutPct\,\% cheaper. Turn-pass fraction moved by \EfQual\ (\EfQualCI), non-inferior, and the
final hidden tests passed in both arms in \EfFinalBoth\ pairs, only at medium in \EfFinalMed\ and only at default in
\EfFinalDef. The preregistered hypothesis was supported. Two sessions in the default arm were flagged and are explained
in the evidence (the flag log): one anchor read \EfSurplusTokens\ more cache tokens than the audit could account
for, worth about half a cent, and one anchor was stopped by the memory guard when the agent's own verification script
ran away. Excluding the flagged pair gives \EfStrict\ (\EfStrictCI). The follow-up cost \$\EfSpend.

\paragraph{What it replaces} The round-2 extrapolation from the Sonnet follow-up is now a measurement. Fable's own
effort effect, $\ln \EfRatio / \ln \RgFabStickyPair$, is about \EfShare\,\% of the log saving of sticky routing on
Fable, not the \RgEffShareFable\,\% extrapolated. Keeping Fable and lowering its effort leaves most of the saving on the
table: exploratory and not concurrent, Fable at medium cost \EfVsSticky\X\ the main campaign's Fable sticky sessions
(\EfVsStickyCI), so Sonnet-at-medium sticky is about \EfStickyVsMedium\X\ of Fable-at-medium, and \EfVsAnchor\X\ the
main campaign's Fable anchors (\EfVsAnchorCI). How Opus responds to effort is still unmeasured.
